\section{Convolutional Neural Network (CNN) Baseline}
\label{sec:cnn_baseline}

The \gls{cnn} architecture \texttt{CNN\_MAX\_REPLICA\_SE} was engineered as the primary baseline to evaluate the \gls{snn} against.
It is a 11.73M parameter \gls{cnn} specifically designed to mirror the \gls{snn} architecture, satisfying the ``Architectural Parity`` functional requirement specified in \ref{subsec:req_functional}. 
This ensures that the only variable being tested in this comparative study is the neural activation mechanism (event-driven binary spikes versus continuous-valued floats), making the comparison immune to common criticisms regarding unmatched or unoptimized baselines.

\subsection{Structural Configuration \& Layer-Wise Mapping}
The \gls{cnn} mirrors the \gls{snn}'s configuration strictly 1:1, replacing every temporal and spiking component with a standard continuous equivalent. Table \ref{tab:snn_cnn_mapping} outlines this direct mapping.

\begin{table}[htbp]
\centering
\resizebox{\textwidth}{!}{%
\begin{tabular}{|l|l|p{8cm}|}
\hline
\textbf{SNN Component} & \textbf{CNN Equivalent} & \textbf{Technical Purpose} \\
\hline
Spatiotemporal Voxel \texttt{[B, 10, 2, H, W]} & Flattened Channel Input \texttt{[B, 20, H, W]} & Translates 10 time bins into a format the respective network can ingest. \\
LIF Neuron (Leaky Integrate-and-Fire) & ReLU Activation & Provides the nonlinearity; SNN uses discrete pulses, \gls{cnn} uses continuous values. \\
Binary Spikes $\{0, 1\}$ & Continuous Floats $[0, \infty)$ & Represents the information ``currency'' being passed between layers. \\
Surrogate Gradient & Exact Derivative (Standard) & Enables backpropagation through the non-differentiable spiking threshold. \\
SEWResBlock (Spiking Shortcut) & Standard Residual Block & Implements identity mapping to prevent vanishing gradients during training. \\
Multi-Step Mode (\texttt{step\_mode=`m'}) & Single spatial pass & Processes the temporal window efficiently on the GPU. \\
Squeeze-and-Excitation (Spike-based) & Squeeze-and-Excitation (Float-based) & Recalibrates channel importance based on global frame context. \\
Temporal Summation ($\sum_{t=1}^{T} S_t$) & N/A (Direct Feature Pass) & Aggregates 100ms of evidence into a single map for the regression head. \\
Multi-Box Head (5 Boxes) & Multi-Box Head (5 Boxes) & Predicts coordinates and objectness scores for 5 targets simultaneously. \\
Synaptic Operations (SOPs) & Floating Point Ops (FLOPs) & The metric used to quantify and compare hardware energy efficiency. \\
\hline
\end{tabular}%
}
\caption{1:1 Architectural Mapping (SNN vs. CNN)}
\label{tab:snn_cnn_mapping}
\end{table}

\subsection{Input Strategy: Channel-Wise Temporal Integration}
Unlike the \gls{snn}, which processes event data sequentially over discrete time steps $T$, the \gls{cnn} baseline utilizes a flattened temporal input strategy. 
Instead of processing 10 timesteps sequentially, it concatenates the $T=10$ bins (each containing 2 polarity channels) into a single input tensor of 20 channels, changing the input shape from $[B, 10, 2, H, W]$ to $[B, 20, H, W]$. 

This approach allows the \gls{cnn} to correlate events across the entire 100ms window in a single spatial pass. This channel-wise stacking of event bins is a standard and effective methodology for feeding asynchronous event-camera data into traditional, non-spiking \gls{cnn} architectures \cite{gehrig2019}.

\subsection{Backbone: Identity Mapping \& Residual Connections}
\label{subsec:cnn_backbone}

The \gls{cnn} backbone employs deep residual learning, utilizing four Residual Blocks that expand to 512 channels in the final stage. 
To satisfy the functional requirement of Architectural Parity, the spiking \texttt{SEWResBlock}s from the \gls{snn} were replaced with Standard Residual Blocks. 
The standard structure within these blocks is $\text{Conv} \rightarrow \text{BatchNorm} \rightarrow \text{ReLU} \rightarrow \text{Conv} \rightarrow \text{BatchNorm} + \text{Shortcut} \rightarrow \text{ReLU}$. 

The inclusion of these skip connections (shortcuts) is vital for enabling the \gls{cnn} network to avoid the same vanishing gradient problem seen in its \gls{snn} counterpart, and successfully learn complex vehicle geometries across the $320 \times 180$ field of view \cite{he2015deep}.

\subsection{SEWResBlock vs Standard ResBlock}
The core mathematical difference between the two architectures lies in their summation logic and the resulting information flow across the shortcut connection:

\begin{itemize}
    \item \textbf{Standard ResBlock (CNN):} 
    Performs an accumulative floating-point addition ($y = f(x) + x$). 
    The continuous intensity of the residual path is directly added to the identity path before the final ReLU activation. 
    Here, the shortcut carries the \textit{raw input features} forward. 
    This mechanism allows the model to learn the ``difference'' (the residual) needed to improve the representation, focusing on information preservation \cite{he2015deep}.
    
    \item \textbf{SEWResBlock (SNN):} 
    Performs a relational element-wise spiking operation \textit{after} the activation function. 
    Because adding two binary spikes (e.g., $1+1=2$) produces a non-binary value, the \gls{snn} uses a spike-merging function (such as ADD or AND) to preserve the discrete $\{0, 1\}$ state. 
    Crucially, the \gls{snn} shortcut carries the \textit{membrane history} forward, prioritizing temporal integrity.

\end{itemize}

Table \ref{tab:resblock_comparison} summarizes these foundational differences in identity mapping.

\begin{table}[htbp]
\centering
\resizebox{\textwidth}{!}{%
\begin{tabular}{|l|p{5cm}|p{5cm}|}
\hline
\textbf{Feature} & \textbf{Standard ResBlock (CNN)} & \textbf{SEWResBlock (SNN)} \\
\hline
Summation & Accumulative ($y = f(x) + x$) & Relational (Spiking Element-Wise) \\
Output Type & Continuous Float $[0, \infty)$ & Binary Spikes $\{0, 1\}$ \\
Primary Benefit & Prevents vanishing gradients & Prevents spike vanishing \\
Mathematical Goal & Information preservation & Temporal integrity \\
\hline
\end{tabular}%
}
\caption[Identity Mapping Comparison]{Comparison of identity mapping between Standard ResBlocks and SEWResBlocks.}
\label{tab:resblock_comparison}
\end{table}

This distinction is particularly relevant for high-speed neuromorphic data: 
while the Standard ResBlock accumulates continuous values that effectively blur fast movements into a single averaged float over the 100ms window, the \gls{snn} equivalent preserves exact spike timing, maintaining the temporal precision required for the regression head.

\subsection{Attention Mechanism: SE-Block}
The \gls{cnn} encounters the same bottleneck as its \gls{snn} counterpart (Section~\ref{subsec:snn_se_block}) and adopts the same \gls{se} mechanism \cite{hu2019_se_block}, differing only in that the squeeze operates over continuous activations via standard global average pooling and the excitation produces a float-valued per-channel scale rather than modulating spike rates. The block confers no hardware-level sparsity benefit in this context; it is retained purely to preserve Architectural Parity with the \gls{snn} baseline (Section~\ref{subsec:req_functional}).

\subsection{The ``Gold Standard'' Upper Bound}
By constructing this custom \gls{cnn} mirror, the baseline serves as the upper theoretical bound of this specific architecture's capability on the ETraM dataset. 
For example, if the \gls{cnn} baseline achieves a Mean \gls{iou} of 48\%, this represents the maximum performance extractable by this parameter configuration. 
If the corresponding SNN achieves 43\% IoU, the 5\% gap explicitly represents the ``quantization cost'' of switching from 32-bit continuous floats to 1-bit binary spikes \cite{cordone2022objectdetectionspikingneural}. 
Adhering to this methodology ensures that the evaluation of the neuromorphic paradigm is both fair and scientifically sound.


\todo[inline,backgroundcolor=myblue,textcolor=white]{Add model architecture diagram.}
\todo[inline,backgroundcolor=myblue,textcolor=white]{Check \& Add model hyperparameter table.}